%!TEX root = ./main.tex

% Styles added for this deck; ran_macros.tex stays identical to the Wi-Fi History copy.

% Navigation symbols sit on top of the source lines at the bottom right.
\setbeamertemplate{navigation symbols}{}

\tikzset{
  % The physical channel: gray, so it reads as "not designed by us".
  chnode/.style={draw=softgray,fill=gray!10,rounded corners=3pt,
    minimum width=1.7cm,minimum height=0.82cm,align=center,font=\bfseries},
  % Shared knowledge / trained model: green, used only for what both ends must hold.
  kbnode/.style={draw=softgreen,fill=green!7,rounded corners=3pt,
    minimum width=1.9cm,minimum height=0.82cm,align=center,font=\bfseries},
  % Learned (neural) blocks: purple, so "trained" is visible at a glance.
  nnnode/.style={draw=ranpurple,fill=purple!6,rounded corners=3pt,
    minimum width=1.9cm,minimum height=0.82cm,align=center,font=\bfseries},
  % Plain content boxes (photo, sentence, answer).
  endnode/.style={draw=softgray,fill=gray!6,rounded corners=3pt,
    minimum width=1.7cm,minimum height=1.0cm,align=center,font=\small\bfseries},
  lbl/.style={font=\scriptsize,text=softgray,align=center},
  iface/.style={font=\scriptsize\bfseries,text=coreorange},
}

% A tiny street-scene video frame (sky, road, one pedestrian, one car), drawn at
% (0,0)--(3.4,2.1). Used on the "lost block" slide.
\tikzset{
  streetframe/.pic={
    \fill[blue!12] (0,1.0) rectangle (3.4,2.1);
    \fill[gray!25] (0,0) rectangle (3.4,1.0);
    \draw[white,thick,dashed] (0,0.45) -- (3.4,0.45);
    \fill[coreorange] (1.15,1.36) circle (0.13);
    \fill[coreorange] (1.03,0.72) rectangle (1.27,1.21);
    \fill[ranblue!70,rounded corners=1pt] (2.25,0.18) rectangle (3.1,0.58);
    \draw[softgray] (0,0) rectangle (3.4,2.1);
  }
}

% Step-by-step reveal for the colored boxes. A tcolorbox inside \uncover would still
% paint its background, so on earlier overlays draw a fully transparent copy instead:
% the box keeps its space and nothing below it jumps when it appears.
% Usage: \takeawayat{3-}{text}   \askat{2-}{text}
\newcommand{\takeawayghost}[1]{%
  \begin{tcolorbox}[colback=orange!7,colframe=coreorange,boxrule=0.7pt,arc=2pt,
    left=5pt,right=5pt,top=4pt,bottom=4pt,
    opacityback=0,opacityframe=0,opacitytext=0]
    \centering\normalsize #1
  \end{tcolorbox}%
}
\newcommand{\askghost}[1]{%
  \begin{tcolorbox}[colback=blue!4,colframe=ranblue,boxrule=0.7pt,arc=2pt,
    left=5pt,right=5pt,top=4pt,bottom=4pt,
    opacityback=0,opacityframe=0,opacitytext=0]
    \centering\large #1
  \end{tcolorbox}%
}
\newcommand{\takeawayat}[2]{\alt<#1>{\takeaway{#2}}{\takeawayghost{#2}}}
\newcommand{\askat}[2]{\alt<#1>{\ask{#2}}{\askghost{#2}}}

% Citation line in the same style as \source, for references with no stable open URL.
\newcommand{\plainsource}[1]{%
  \par\vspace{0.25em}%
  {\tiny\textcolor{gray}{#1}\par}%
}

% Picture-based "signal gets weaker" sketches (slides "Where it struggles" and
% "Weak signal: blurrier, not frozen"). All frames use the streetframe geometry.
% \signalbars{x}{y}{n}: four bars at (x,y), the first n filled.
\newcommand{\signalbars}[3]{%
  \foreach \i in {1,2,3,4}{%
    \ifnum\i>#3\relax
      \draw[softgray,thick] (#1+0.22*\i,#2) rectangle ++(0.14,0.12*\i);
    \else
      \fill[ranblue] (#1+0.22*\i,#2) rectangle ++(0.14,0.12*\i);
    \fi}%
}
% Pedestrian and car at a given opacity, so blurred frames can stack shifted copies.
\newcommand{\streetobjects}[1]{%
  \fill[coreorange,opacity=#1] (1.15,1.36) circle (0.13);
  \fill[coreorange,opacity=#1] (1.03,0.72) rectangle (1.27,1.21);
  \fill[ranblue!70,opacity=#1,rounded corners=1pt] (2.25,0.18) rectangle (3.1,0.58);%
}
\tikzset{
  % today's link below its threshold: the picture stops
  frozen/.pic={
    \fill[black,opacity=0.6] (0,0) rectangle (3.4,2.1);
    \fill[white] (1.45,0.85) rectangle (1.6,1.45);
    \fill[white] (1.8,0.85) rectangle (1.95,1.45);
    \node[font=\small\bfseries,text=white] at (1.7,0.45) {frozen};
  },
  % today's link after feedback lowered the bitrate: coarser, blockier picture
  streetblocky/.pic={
    \fill[blue!12] (0,1.0) rectangle (3.4,2.1);
    \fill[gray!25] (0,0) rectangle (3.4,1.0);
    \fill[coreorange!85] (0.98,1.2) rectangle (1.32,1.54);
    \fill[coreorange!85] (0.98,0.68) rectangle (1.32,1.2);
    \fill[ranblue!60] (2.2,0.12) rectangle (3.15,0.64);
    \draw[white,opacity=0.55] (0,0) grid[step=0.34] (3.4,2.1);
    \draw[softgray] (0,0) rectangle (3.4,2.1);
  },
  % today's link right after a sudden drop, before feedback reacts: frames arrive late
  late/.pic={
    \fill[black,opacity=0.45] (0,0) rectangle (3.4,2.1);
    \draw[white,very thick] (1.7,1.25) circle (0.38);
    \draw[white,very thick] (1.7,1.25) -- (1.7,1.5);
    \draw[white,very thick] (1.7,1.25) -- (1.88,1.25);
    \node[font=\small\bfseries,text=white] at (1.7,0.45) {late};
  },
  % Version 3: coarse background, sharp pedestrian (bits spent where the task looks)
  streetv3/.pic={
    \fill[blue!12] (0,1.0) rectangle (3.4,2.1);
    \fill[gray!25] (0,0) rectangle (3.4,1.0);
    \fill[ranblue!60] (2.2,0.12) rectangle (3.15,0.64);
    \draw[white,opacity=0.55] (0,0) grid[step=0.34] (3.4,2.1);
    \fill[blue!12] (0.82,1.0) rectangle (1.48,1.68);
    \fill[gray!25] (0.82,0.58) rectangle (1.48,1.0);
    \fill[coreorange] (1.15,1.36) circle (0.13);
    \fill[coreorange] (1.03,0.72) rectangle (1.27,1.21);
    \draw[softgreen,thick] (0.82,0.58) rectangle (1.48,1.68);
    \draw[softgray] (0,0) rectangle (3.4,2.1);
  },
  % learned link with a weaker signal: same scene, softer
  streetsoft/.pic={
    \fill[blue!12] (0,1.0) rectangle (3.4,2.1);
    \fill[gray!25] (0,0) rectangle (3.4,1.0);
    \draw[white,thick,dashed,opacity=0.6] (0,0.45) -- (3.4,0.45);
    \foreach \dx/\dy in {-0.05/0,0.05/0,0/0.05,0/-0.05}{
      \begin{scope}[shift={(\dx,\dy)}]\streetobjects{0.35}\end{scope}}
    \draw[softgray] (0,0) rectangle (3.4,2.1);
  },
  streetblurry/.pic={
    \fill[blue!12] (0,1.0) rectangle (3.4,2.1);
    \fill[gray!25] (0,0) rectangle (3.4,1.0);
    \draw[white,thick,dashed,opacity=0.35] (0,0.45) -- (3.4,0.45);
    \foreach \dx/\dy in {-0.12/0,0.12/0,0/0.12,0/-0.12,-0.08/-0.08,0.08/0.08,
        -0.08/0.08,0.08/-0.08}{
      \begin{scope}[shift={(\dx,\dy)}]\streetobjects{0.2}\end{scope}}
    \draw[softgray] (0,0) rectangle (3.4,2.1);
  },
}
